% ======================================================================
% sec-prelim.tex -- the determination model (instantiation-agnostic).
%
% Self-contained, but notation and definitions deliberately match a
% determination framework: specifications are relations with input/output
% refinement orders; commitments are fateful restrictions of a residual
% specification; exposures/entailments are content-preserving free
% transitions; determinations are first-hit complete layer traces.
% ======================================================================
\section{The Determination Model}
\label{sec:prelim}

We formalize the setting in which provenance questions become ambiguous:
specifications that admit several admissible outcomes per input, and the
irrevocable choices---\emph{commitments}---that narrow this ambiguity until a
single outcome is determined. The development is self-contained; it isolates
exactly the structure the provenance algebra of
Section~\ref{sec:algebra} requires.

\subsection{Specifications}

A \emph{specification} relates inputs to the outputs considered acceptable
for them. Both inputs and outputs carry a \emph{refinement order}: a partial
order under which a larger element adds detail without contradicting a
smaller one.

\begin{definition}[Specification]
  \label{def:spec}
  A \emph{specification} is a tuple
  $\Spec = (I, O, \OrdI, \OrdO, S)$ where $I$ is a set of inputs ordered by
  $\OrdI$, $O$ is a set of outputs ordered by $\OrdO$, and
  $S \subseteq I \times O$ is a relation assigning to each input at least one
  acceptable output. We write $S(i) \triangleq \{o \in O \mid (i,o) \in S\}$
  for the \emph{acceptable set} at $i$.
\end{definition}

\noindent
The output order encodes which replacements are safe: when
$o_1 \OrdO o_2$, replacing $o_1$ by $o_2$ \emph{refines} it (adds detail)
rather than contradicting it. The order is part of the specification---a
modeling choice, not something read off execution structure. Enlarging
$\OrdO$ makes more outputs mutually compatible and so reduces the commitment
needed to resolve ambiguity. When the input order is equality we recover an
ordinary fixed-input problem; this paper's instantiations fix the input, so
we write $\Spec$ for the fixed-input restriction and suppress $\OrdI$ except
where it is needed.

A specification is in general a \emph{relation}; the goal of determination is
to turn it into a \emph{function}---one determined output per input. Because
provenance analysis is retrospective, we read the determined output as the
$\OrdO$-maximum of the acceptable set.

\begin{definition}[Determined]
  \label{def:determined}
  $\Spec$ is \emph{determined} at $i$ if $S(i)$ is a chain under $\OrdO$ with
  a maximum element; the determined output is $\Max_{\OrdO} S(i)$. $\Spec$ is
  \emph{determined} if it is determined at every input. Ambiguity requiring
  resolution arises exactly when $S(i)$ contains $\OrdO$-incomparable outputs.
\end{definition}

\noindent
A determined specification is a function from inputs to outputs; classical
provenance applies to it directly. Everything new in this paper concerns what
happens \emph{before} determination, when incomparable outputs coexist.

\subsection{Residual States and Commitments}

Resolution proceeds by restricting the acceptable relation. We track its
progress as a \emph{semantic state}.

\begin{definition}[Semantic state]
  \label{def:state}
  A \emph{semantic state} over $\Spec$ on a finite input set $J \subseteq I$
  is a pair $\sigma = (J, S_\sigma)$ with $S_\sigma \subseteq S|_J$ a residual
  subrelation in which every input of $J$ retains at least one acceptable
  output. The \emph{initial state} is $\sigma_0 = (J, S|_J)$. A state is
  \emph{determined} when each input's residual fiber is a singleton.
\end{definition}

\noindent
Resolution only ever \emph{removes} rows from the residual; it never adds
them. Two kinds of event drive this process, and the distinction between them
is the crux of the model.

\begin{definition}[Commitment and basis]
  \label{def:commitment}
  A \emph{commitment} $\varphi$ is a labeled partial operator on semantic
  states: where it is defined at $\sigma$ it leaves the input set unchanged
  and deletes at least one row from the residual,
  $S_{\varphi(\sigma)} \subsetneq S_\sigma$, with $S_{\varphi(\sigma)}$
  nonempty. A commitment is \emph{irrevocable}: once applied, the deleted rows
  do not return. A \emph{commitment basis} $\Basis$ is the set of commitments
  available for a given specification; it is part of the application model
  (the interface of binding acts the system can take), not merely a
  mathematical generating set.
\end{definition}

\noindent
The basis is a modeling choice: different bases yield different determination
structures and so different provenance. A commitment
is \emph{fateful}: some output that was acceptable before it is
no longer attainable afterwards. We contrast this with changes that remove
rows while preserving what is attainable.

\begin{definition}[Content-preserving restriction; free transitions]
  \label{def:free}
  A state change from $\sigma$ to $\tau$ (with $S_\tau \subseteq S_\sigma$) is
  \emph{content-preserving} if every removed row has, at the same input, a
  retained row refining it under $\OrdO$; otherwise it is \emph{fateful}.
  Commitments are fateful. An \emph{input update} (the environment supplying a
  larger input, removing rows dominated under $\OrdI$) and an \emph{exposure}
  (reporting a determined-enough output, narrowing a fiber content-preservingly)
  are content-preserving \emph{free transitions}: they change the residual but
  make no commitment, and---crucially for what follows---carry no work and no
  depth. An \emph{entailment} is the output-side counterpart: it refines the
  observed output upward in $\OrdO$ without excluding any admissible
  alternative.
\end{definition}

\noindent
Separating commitments (which exclude an alternative outright) from exposures
and entailments (which only add compatible detail) is what lets us count
\emph{resolution} without conflating it with reporting or with the arrival of
new input. In our lead instantiation the free transitions are concrete: in
$\mbox{Datalog}^\neg$, deriving a positive consequence of a sealed stratum is
an entailment, whereas committing a negative-cycle atom to a truth value is a
commitment.

\subsection{Determinations}

Several commitments typically commute; grouping them records that they could
be issued together without ordering among them.

\begin{definition}[Parallel step; layer]
  \label{def:parallel-step}
  A \emph{parallel step} at $\sigma$ is a finite nonempty set
  $L \subseteq \Basis$ of commitments, each applicable at the same source
  $\sigma$, together with a declared semantic state $\tau$ whose residual is
  their common-source conjunction,
  $S_\tau = \bigcap_{\varphi\in L} S_{\varphi(\sigma)}$; we write
  $\sigma \xrightarrow{L} \tau$. The recorded step occupies one
  \emph{layer} of resolution. Validity is a property of the whole set $L$ at
  its source: pairwise commutation alone need not suffice in the general
  model, and a recorded step is not retrospectively enlarged to a maximal
  one.
\end{definition}

\begin{definition}[Determination; work and depth]
  \label{def:determination}
  A \emph{determination} from $\sigma_0$ over $\Basis$ is a trace of parallel
  steps
  \[
    D:\quad
    \sigma_0 \xrightarrow{L_1} \sigma_1 \xrightarrow{L_2} \cdots
    \xrightarrow{L_m} \sigma_m .
  \]
  It is \emph{complete} (first-hit) if $\sigma_m$ is the first determined
  state on the trace; a complete determination is also called a
  \emph{resolving} determination, as it resolves the specification to a
  single output. Its \emph{work} and \emph{depth} are
  \[
    \operatorname{work}(D) = \sum_{j=1}^{m} |L_j|,
    \qquad
    \operatorname{depth}(D) = m .
  \]
  Write $\Dets = \Dets_{\sigma_0,\Basis}$ for the space of complete
  determinations. Free transitions are not parallel steps and are never
  coalesced into a layer; an online trace is analyzed segment by segment
  between its free transitions.
\end{definition}

\noindent
Each concrete execution of a real system corresponds to one determination---the
commitments it actually made. The complete determinations are prefix-free:
no complete determination is a layer-prefix of another, since a trace stops at
the \emph{first} determined state. Distinct determinations may reach the same
determined result; the \emph{result map}
\[
  \rho : \Dets \longrightarrow \Rho(J)
\]
sends each complete determination to the determined (functional) result it
produces, where $\Rho(J)$ is the set of functional relations on $J$ contained
in $S$. The result records \emph{what} was determined, not \emph{which} trace
led there.

\begin{remark}[Why layers, not single commitments]
  \label{rem:layers}
  Commitments are grouped into layers for two reasons. First, non-commitment
  events may intervene, so a later commitment may depend on an intervening
  event and cannot be reordered before it. Second, even contiguous commitments
  need not commute when one's effect depends on the other's exclusions. The
  constitutional notion is the \emph{actual valid common-source step} recorded
  by the trace, not a maximal pairwise-commuting set. In both instantiations of
  this paper, however, the declared commitment families are
  state-independently independent, so their valid common-source conjunctions
  are exactly the commuting batches used below; this is an instantiation
  property, not part of the general definition.
\end{remark}
